\subsection{Hypothèses communes à toutes les résolutions}
L'énoncé nous dit de considérer que les frottements de l'air sont négligeables. 
Donc l'énergie mécanique d'un corps va être conservée, et pourra toujours s'écrire sous la forme:
\begin{gather*}
    \text{Energie mécanique} = \text{énergie potentielle} + \text{énergie cinétique} \\
    E_m = E_c + E_p
\end{gather*}
où l'énergie mécanique reste constante, et l'énergie potentielle va descendre pendant que l'énergie cinétique va monter, ou le contraire (transvasements entre énergie potentielle et cinétique, sans aucune perte).

\paragraph{Parenthèse sur les frottements}
Dans le cas du saut à l'élastique où les vitesses maximales sont faibles et où les personnes qui sautent n'ont pas de parachute, les frottements de l'air sont très probablement négligeables. 

Dans d'autres situations, ce ne l'est pas.

On se souvient tous en effet de l'expérience d'une plume et d'une pierre qu'on laisse tomber dans le vide qui touchent le sol au même instant, alors que dans l'air, la plume descend beaucoup plus lentement que la pierre.

Ou encore lors de la rentrée d'une capsule spatiale dans l'atmosphère, on voit que l'énergie mécanique est convertie en chaleur à cause des frottements de l'air qui sont si importants à cette vitesse que la capsule devient rouge comme une braise